\par\SourceLeaf{134}
Gerade dadurch erhält nun aber die Unklarheit\footnote{Wer an dieser Bezeichnung einer seit einem halben Jahrhundert geltenden Anordnung Anstoß nimmt, möge sein Urteil zurückstellen, bis er die erste Abteilung ganz gelesen hat.} des Regulativs bei derjenigen Bestimmung, durch die die Höhe und das Fortschreiten des Pensionssatzes geregelt werden sollen, einen sehr bedenklichen Charakter. Die eigenen Worte des Reskripts des Universitätskuratoriums vom 20. November 1794, durch das diese Regelung genehmigt wurde, lauten:
\begin{quote}
— — — — Zweitens genehmigen wir, dass jedes Mal, wenn sich das Vermögen der Kasse um 5000 Thl. vermehrt hat und die anderen Einnahmen des Witwenfonds keine Verminderung erfahren haben und auch die Anzahl der Pensionen nicht über 15 gestiegen ist, jede Pension um 10 Thl. erhöht werden soll.
\end{quote}
\noindent\hspace*{\parindent}Mangelhaft ist diese Bestimmung darin, dass sich nicht ganz unzweideutig erkennen lässt, \textit{was} an einer Pensionserhöhung eigentlich von den vorangehenden Bedingungen abhängig sein soll: ihr Fortbestehen oder bloß ihr Anfang; mit anderen Worten (wobei ich nur die letzte Bedingung betrachte),
\begin{quote}
ob eine Erhöhung, die dem Statut gemäß zu einer Zeit eingetreten ist, als die Witwenzahl höchstens 15 betrug, wieder aufhören oder wenigstens zweckmäßig verändert werden soll, sobald später die Anzahl der Pensionen jene festgelegte Zahl überschreitet,
\end{quote}
oder aber,
\begin{quote}
ob eine einmal eingetretene Erhöhung unabhängig von der späteren Überschreitung dieser Zahl dennoch unabänderlich fortdauern soll.
\end{quote}
\noindent\hspace*{\parindent}Dass die obige Formulierung der Vorschrift natürlicher zur zweiten Auslegung führt als zur ersten, wenn man ohne jede Rücksicht auf die Folgerungen aus ihrem Inhalt nur den Wortlaut betrachtet, will ich umso weniger bestreiten, als der sprachliche Gesichtspunkt bei meiner gegenwärtigen, hauptsächlich auf den inneren Gehalt der Anordnung selbst gerichteten Untersuchung ganz untergeordnet ist. Ich kann jedoch nicht umhin, zum Vergleich auch die Fassung hierherzusetzen, in der \textsc{Brandes} die Verfügung in seinem bekannten Werk über die Universität Göttingen anführt. Da \textsc{Brandes} 1794 im Ministerium als Referent für die Universitätsangelegenheiten tätig war, gehört seine Auffassung jedenfalls zur Sache, auch wenn sie vom juristischen Gesichtspunkt aus, der meiner Untersuchung ebenfalls fremd ist, die im offiziellen Reskript einmal gebrauchte Wortfassung nicht außer Kraft setzen kann. An der angegebenen Stelle, S. 254, heißt es:
\begin{quote}
Ferner wurde damals beschlossen, dass jedes Mal, wenn das Kapitalvermögen der Kasse um 5000 Thl. angewachsen sei, jede Pension um 10 Thl. vermehrt werden solle, \textit{solange} die Zahl der Pensionierten nicht über 15 hinausgeht, was noch nie der Fall war.
\end{quote}
\noindent\hspace*{\parindent}Danach scheint \textsc{Brandes} den Vorbehalt hinsichtlich der 15 Pensionen eher im Sinn der \textit{ersten} Interpretation verstanden, ihn aber wegen des zuletzt von ihm angeführten Umstands für praktisch ganz unerheblich gehalten zu haben. Dadurch lässt sich vielleicht auch erklären, dass bei der Wortfassung des Reskripts nicht für die größtmögliche Genauigkeit gesorgt wurde.

\textsc{Brandes} hätte jedoch wahrscheinlich ganz anders geurteilt, wenn er vorher die Tatsachen genau geprüft hätte. \textsc{Brandes} kann mit der Angabe »was noch nie der Fall war« nicht den siebenjährigen Zeitraum von der Einführung der Bestimmung bis zur Abfassung seines Buches gemeint haben, da der Erfahrung aus einem so kurzen Zeitraum überhaupt kein Gewicht beigelegt werden könnte. Er muss vielmehr die gesamte seit der Stiftung der Kasse vergangene Zeit gemeint haben. Dann ist seine Behauptung aber sachlich unrichtig. Die Zahl der aus der Kasse bezahlten Pensionen war tatsächlich schon einmal über 15 gestiegen, nämlich während der ersten sechs Monate des Jahres 1778, als 16 Pensionen bestanden; derselbe Fall wäre sogar schon
\SourceLeaf{135}
zwei Jahre früher eingetreten, wenn die erst 1794 eingeführte Verlängerung der Waisenpensionen bis zum Alter von 20 Jahren schon damals gegolten hätte. Dieser Umstand erhält jedoch ein sehr ernstes Gewicht, wenn man gemäß der oben gemachten Bemerkung bedenkt, dass 20—30 Jahre vor jener Zeit, nämlich von 1747—1757, die durchschnittliche Teilnehmerzahl nur 21—22 betrug, 1794—1802 dagegen 35—36.

Über 15 ist die Zahl der Pensionen danach erst wieder Ostern 1837 gestiegen; seitdem ist sie stets darüber geblieben. Die Kasse hat den vollen Pensionsbetrag einschließlich aller seit 1794 eingetretenen aufeinanderfolgenden Erhöhungen weitergezahlt und damit wenigstens \textit{implizit} — denn ob die Kirchendeputation vorher darüber verhandelt hat, ist mir nicht bekannt — die zweite Interpretation des zweifelhaften Punkts angenommen.

Wollen wir nun aber unter Beiseitelassung sprachlicher und formell juristischer Rücksichten zwischen den beiden Interpretationen nur nach ihren inneren Gründen entscheiden, so drängt sich zunächst sofort die Frage auf: Wenn es wirklich unbedenklich ist, eine erhöhte Pension auch für 16 Beziehende ungeschmälert \textit{fortbestehen} zu lassen, warum soll es dann verboten sein, sie auch bei 16 Beziehenden \textit{anfangen} zu lassen? Die Gefahr, wenn eine besteht, ist im einen Fall doch genau so groß wie im anderen. Dieses Argument wird noch schlagender, wenn man zu größeren Zahlen übergeht. Denn eine erhöhte Pension bei 17 Beziehenden (und wie viel mehr bei 18, 19 und so weiter) fortbestehen zu lassen, ist doch ganz offenbar für die Kasse gefährlicher, als die Erhöhung bei 16 anfangen zu lassen. Nimmt man die zweite Interpretation an, so ist aber jenes erlaubt und dieses verboten.

Zu welchen Folgerungen eine streng konsequente Durchführung der zweiten Interpretation führt, lässt sich leicht überblicken.

Im gedruckten Regulativ von 1838 wird die fragliche Anordnung in §. 10 mit den Worten eingeleitet: Die Witwenpension \textit{wird} teils nach dem Vermögensstand der Kasse, teils nach der Zahl der Witwen \textit{bestimmt}. Nach der Einleitung zur ersten Veröffentlichung des Regulativs (1833) soll dieses die Verpflichtungen und Rechte der Teilnehmer feststellen. Ein neues Mitglied tritt bei, indem es ein ihm zur Kenntnisnahme der Pflichten und Rechte vorgelegtes Exemplar unterschreibt. Man muss daher annehmen, dass das Regulativ diese \textit{vollständig} enthält und dass die Kasse gegenüber den Mitgliedern oder deren Hinterbliebenen keine Rechte geltend machen kann, die nicht in diesem Regulativ enthalten sind.

Nun findet sich darin aber kein einziges Wort, das ein Recht vorbehält, den Pensionssatz gegebenenfalls wieder herabzusetzen, es sei denn, man nimmt die erste Interpretation jener Bestimmung des Pensionssatzes nach der Witwenzahl an. Anderenfalls muss man bei strenger Konsequenz einräumen, dass die Kasse verpflichtet sei, sämtliche Erhöhungen, die bei einer Pensionszahl unter 16 eingetreten sind, ungeschmälert weiterzuzahlen, die Zahl der Witwen (und Waisen) möge später (zum Beispiel infolge des insgesamt erweiterten Teilnehmerkreises) \textit{so hoch anwachsen, wie sie wolle.}

Das ist aber geradezu ungereimt, sofern man nicht annimmt, dass die Regierung die Gewähr leisten müsse; dies zu beurteilen liegt außerhalb meiner Zuständigkeit.

